\section{符号说明}
\label{sec:符号说明}

本书采用如下字型约定：\textbf{向量与矩阵一律用粗体斜体（$\boldsymbol{}$）表示}，例如 $\boldsymbol{a}$、$\boldsymbol{A}$；\textbf{标量、分量与矩阵元素不加粗}，例如 $a_i$、$a_{ij}$。单位向量组用 $\mathbf{}$（正粗体）表示，如 $\mathbf{a}_i$、$\mathbf{b}_i$。凡表示"在参考系 $x$ 中表达或取分量"的下标一律写成 $|x$，如 $\nu_{i|A}$。相对转动量用竖线下标，如 $\boldsymbol{\omega}_{B|A}$ 表示 $B$ 在 $A$ 中的角速度。

此处列出本课程所用的符号，可按下列格式分表整理（可参照《符号约定与工作要求.md》）：

\begin{table}[H]
\centering
\caption{符号说明（示例）}
\label{tab:符号说明}
\begin{tabular}{ll}
\toprule
符号 & 含义 \\
\midrule
$A$,\, $B$ & 参考系（或刚体） \\
$\boldsymbol{a}$ & 向量 \\
$a_{ij}$ & 矩阵元素（标量） \\
\bottomrule
\end{tabular}
\end{table}
